\documentclass[10pt,a4paper]{amsart}
\usepackage[margin=19mm]{geometry}
\usepackage{amsmath,amssymb,mathtools}
\usepackage[scr=rsfs,cal=euler]{mathalfa}
\usepackage{xfrac}
\usepackage[unicode,hidelinks,bookmarks=false]{hyperref}
\newcommand{\ie}{i.e.\ }
\newcommand{\cf}{cf.\ }
\newcommand{\resp}{respectively\ }
\newcommand{\Subsection}{\subsection}
\providecommand{\nobreakdash}{\nobreak\hbox{-}}
\setlength{\emergencystretch}{2em}
\multlinegap0pt
\usepackage{mathtools}
\usepackage{amssymb}
\usepackage[shortlabels]{enumitem}
\usepackage{csquotes}
\usepackage{bm}
\newtheorem{lemma}{Lemma}
\DeclareMathOperator{\N}{\mathbb{N}}
\DeclareMathOperator{\R}{\mathbb{R}}
\newcommand{\ccint}[1]{\left[{#1}\right]} 
\newcommand{\coint}[1]{\left[{#1}\right)} 
\newcommand{\modu}[1]{\left\vert{#1}\right\vert}
\title{The $p$-spectrum of Random Wavelet Series}
\author{C\'eline Esser$^a$ , Thelma Lambert$^a$ and B\'eatrice Vedel$^b$}
\date{Source: arXiv:2510.00622v2 (CC BY 4.0)}
\begin{document}
\maketitle

\noindent\emph{Source and licence.} The $p$-spectrum of Random Wavelet Series. Version arXiv:2510.00622v2 (CC BY 4.0), distributed by arXiv under the Creative Commons Attribution 4.0 International licence, http://creativecommons.org/licenses/by/4.0/, as recorded in the arXiv licence metadata for this exact version and shown on the abstract page https://arxiv.org/abs/2510.00622v2. Full licence notice: \texttt{licenses/arxiv-2510.00622-CC-BY-4.0.txt}.\par\smallskip

\noindent\textit{Editorial scope.} Counted unit: the article's lemma lm:rid\_order2 (source0.tex lines 877--890). The article's complete original proof of it is delivered with it (source0.tex lines 892--987, one \texttt{proof} environment of 96 lines). No mathematics is written, repaired or re-derived: the statement and the proof are the article's own lines, in the article's own notation, and no retained line is substituted or re-flowed.  Retained uncounted as the source-local context the proof needs: lines 200--203; lines 636--638; lines 709--711; lines 718--724; lines 764--766; lines 768--770.  Nothing was omitted from any retained span, so the package carries no omission record.  No retained line contains a citation, so the wrapper carries no bibliography and no bibliography entry.  No retained line contains a figure environment, an image-inclusion command or drawing code, so no image is delivered and the package declares no asset dependency.  Editorial formatting only, in addition: an AMS \texttt{amsart} preamble that restates the article's own declarations and macros used by the retained lines, namely the environments \texttt{lemma} on separate counters, together with the standard packages those lines need.  The retained environments print in this document as Lemma 1 in order of appearance, whereas the article prints the same items with the numbering of its own full document.  The complete native source of the article as distributed in the arXiv e-print for this exact version is bundled unchanged as \texttt{sources/186148/source0.tex}.
\begin{equation}\label{eq:bound_rho}
    \rho_{\vec{c}}(\alpha) \leq \alpha p+1 -\eta_f(p)
\qquad \text{for every }\alpha\in\R,
\end{equation}

\begin{equation}\label{eq:epsilon_neg}
    h+2\varepsilon ,\; M\varepsilon p +\frac{Mhp}{2}+1<0 \quad\text{if }h<0, 
\end{equation}

\begin{equation}\label{eq:delta}
    0<\delta<\min\left(\frac{1}{5m},\frac{1}{5}\,\rho_{\vec{c}}^{(p),\ast}(h),\eta_f(p),\frac{\rho^{(p),\ast}_{\vec{c}}(h)-\displaystyle\sup_{\alpha\in\coint{h,0}}\left(\frac{(\rho_{\vec{c}}(\alpha)-1)h}{\alpha}+1\right)-M\varepsilon p}{4+\frac{h\left(h+2\varepsilon+\frac{1}{p}\right)}{\left(\alpha_0+\frac{1}{p}\right)(h+2\varepsilon)}}\right) 
\end{equation}

\begin{lemma}\label{lm:card_Lambdap}
    There exists $A\in\mathcal{A}$ and a sequence $(j_n')_{n\in\N}$ such that for every $n\in\N$ large enough,
    \[
        \#\widetilde{\Lambda}_{j_n}^{(p)}(h,\varepsilon) \coloneqq \#\left\{\lambda\in\Lambda_{j_n}^{(p)}(h,\varepsilon):A(\lambda)=A \text{ and }j_n'(\lambda)=j_n'\right\} \geq 2^{\left(\rho_{\vec{c}}^{(p),\ast}(h)-3\delta\right)j_n}.
    \]
    In particular, for every $n\in\N$, $j_n'\in\ccint{Aj_n,\left(A+\frac{1}{m}\right)j_n}$.
\end{lemma}

\begin{equation}\label{eq:epsilon_s1}
    \varepsilon(\alpha)<\frac{\varepsilon}{6\left(A+\frac{1}{m}\right)}
\end{equation}

\begin{equation}\label{eq:epsilon_s2}
    \#\left\{\lambda'\in\Lambda_{j'}:2^{-\left(\alpha-\frac{1}{p}+\varepsilon(\alpha)\right)j'} \leq \modu{c_{\lambda'}} \leq 2^{-\left(\alpha-\frac{1}{p}-\varepsilon(\alpha)\right)j'}\right\} \leq 2^{\left(\rho_{\vec{c}}\left(\alpha-\frac{1}{p}\right)+\delta\right)j'}.
\end{equation}

\begin{lemma}\label{lm:rid_order2}
    If $h<0$, then for every $s\in\{1,\ldots,S\}$ such that $A\geq\frac{h+2\varepsilon}{\alpha_s-\frac{1}{p}}$, for every $n\in\N$ large enough, if $\widetilde{\widetilde{\Lambda}}_{j_n}^{(p)}(h,\varepsilon)$ denotes the set of dyadic intervals $\lambda\in\widetilde{\Lambda}_{j_n}^{(p)}(h,\varepsilon)$ satisfying
    \[
    \#\left\{\lambda'\in\Lambda_{j_n'}:\lambda'\subseteq\lambda  \text{ and } 2^{-\left(\alpha_s-\frac{1}{p}+\varepsilon_s\right)j_n'} \leq \modu{c_{\lambda'}} \leq 2^{-\left(\alpha_s-\frac{1}{p}-\varepsilon_s\right)j_n'}\right\}<2^{-\left(h+\frac{11\varepsilon}{6}+\frac{1}{p}\right)pj_n}2^{\alpha_spj_n'},
    \]
    then
    \[
    \#\widetilde{\widetilde{\Lambda}}_{j_n}^{(p)}(h,\varepsilon) \geq 2^{\left(\rho_{\vec{c}}^{(p),\ast}(h)-4\delta\right)j_n}
    \]
    and for every $\lambda\in\widetilde{\widetilde{\Lambda}}_{j_n}^{(p)}(h,\varepsilon)$, one has
    \[
    \sum_{\lambda'\subseteq\lambda \,:\, 2^{-\left(\alpha_s-\frac{1}{p}+\varepsilon_s\right)j_n'} \leq \modu{c_{\lambda'}} \leq 2^{-\left(\alpha_s-\frac{1}{p}-\varepsilon_s\right)j_n'}} \modu{c_{\lambda'}}^p2^{-(j_n'-j_n)} \leq 2^{-\left(h+\frac{5\varepsilon}{3}\right)pj_n}.
    \]
\end{lemma}

\begin{proof}
    First, for all $\lambda\in \widetilde{\widetilde{\Lambda}}_{j_n}^{(p)}(h,\varepsilon)$, 
    \begin{align*}
        & \sum_{\lambda'\subseteq\lambda \,:\, 2^{-\left(\alpha_s-\frac{1}{p}+\varepsilon_s\right)j_n'} \leq \modu{c_{\lambda'}} \leq 2^{-\left(\alpha_s-\frac{1}{p}-\varepsilon_s\right)j_n'}} \modu{c_{\lambda'}}^p2^{-(j_n'-j_n)} \\
        & \leq \#\left\{\lambda'\subseteq\lambda :2^{-\left(\alpha_s-\frac{1}{p}+\varepsilon_s\right)j_n'} \leq \modu{c_{\lambda'}} \leq 2^{-\left(\alpha_s-\frac{1}{p}-\varepsilon_s\right)j_n'}\right\} 2^{-\left(\alpha_s-\frac{1}{p}-\varepsilon_s\right)pj_n'}2^{-(j_n'-j_n)} \\
        & \leq 2^{-\left(h+\frac{11\varepsilon}{6}+\frac{1}{p}\right)pj_n}2^{\alpha_spj_n'}2^{-(\alpha_s-\varepsilon_s)pj_n'}2^{j_n} \\
        & \leq 2^{-\left(h+\frac{11\varepsilon}{6}\right)pj_n}2^{\varepsilon_spj_n'} \\
        & \leq 2^{-\left(h+\frac{5\varepsilon}{3}\right)pj_n},
    \end{align*}
    using Condition~\eqref{eq:epsilon_s1}.

    Secondly, let us assume on the contrary that 
    \[
    \#\widetilde{\widetilde{\Lambda}}_{j_n}^{(p)}(h,\varepsilon) < 2^{\left(\rho_{\vec{c}}^{(p),\ast}(h)-4\delta\right)j_n}.
    \]
    Then, using Equation~\eqref{eq:epsilon_s2} and Lemma~\ref{lm:card_Lambdap},
    \begin{align*}
        2^{\left(\rho_{\vec{c}}\left(\alpha_s-\frac{1}{p}\right)+\delta\right)j_n'} 
        & \geq \#\left\{\lambda'\in\Lambda_{j_n'} :2^{-\left(\alpha_s-\frac{1}{p}+\varepsilon_s\right)j_n'} \leq \modu{c_{\lambda'}} \leq 2^{-\left(\alpha_s-\frac{1}{p}-\varepsilon_s\right)j_n'}\right\} \\
        & \geq \sum_{\lambda\in\widetilde{\Lambda}_{j_n}^{(p)}(h,\varepsilon)\setminus\widetilde{\widetilde{\Lambda}}_{j_n}^{(p)}(h,\varepsilon)}\#\left\{\lambda'\in\Lambda_{j_n'} :\lambda'\subseteq\lambda \text{ and } 2^{-\left(\alpha_s-\frac{1}{p}+\varepsilon_s\right)j_n'} \leq \modu{c_{\lambda'}} \leq 2^{-\left(\alpha_s-\frac{1}{p}-\varepsilon_s\right)j_n'}\right\} \\
        & \geq \#\left(\widetilde{\Lambda}_{j_n}^{(p)}(h,\varepsilon)\setminus\widetilde{\widetilde{\Lambda}}_{j_n}^{(p)}(h,\varepsilon)\right)2^{-\left(h+\frac{11\varepsilon}{6}+\frac{1}{p}\right)pj_n}2^{\alpha_spj_n'} \\
        & > \left(2^{\left(\rho_{\vec{c}}^{(p),\ast}(h)-3\delta\right)j_n}-2^{\left(\rho_{\vec{c}}^{(p),\ast}(h)-4\delta\right)j_n}\right)2^{-\left(h+2\varepsilon+\frac{1}{p}\right)pj_n}2^{\alpha_spj_n'} \\
        & = 2^{\left(\rho_{\vec{c}}^{(p),\ast}(h)-3\delta\right)j_n}\left(1-2^{-\delta j_n}\right)2^{-\left(h+2\varepsilon+\frac{1}{p}\right)pj_n}2^{\alpha_spj_n'} \\
        & \geq 2^{\left(\rho_{\vec{c}}^{(p),\ast}(h)-3\delta\right)j_n}2^{-\delta j_n}2^{-\left(h+2\varepsilon+\frac{1}{p}\right)pj_n}2^{\alpha_spj_n'},
    \end{align*}
    i.e.
    \[
    2^{\left(\rho_{\vec{c}}\left(\alpha_s-\frac{1}{p}\right)+\delta-\alpha_sp\right)j_n'} \geq 2^{\left(\rho_{\vec{c}}^{(p),\ast}(h)-4\delta-(hp+1)-2\varepsilon p\right)j_n}.
    \]
    However, this inequality cannot hold. Indeed, by Equation~\eqref{eq:bound_rho}
and the condition $\delta<\eta_f(p)$, we have
\[
\rho_{\vec{c}}\left(\alpha_s-\frac{1}{p}\right)
+\delta-\alpha_sp<0.
\]
Since $j_n'\geq Aj_n$, it follows that
\[
2^{\left(\rho_{\vec{c}}\left(\alpha_s-\frac{1}{p}\right)
+\delta-\alpha_sp\right)j_n'}
\leq
2^{\left(\rho_{\vec{c}}\left(\alpha_s-\frac{1}{p}\right)
+\delta-\alpha_sp\right)Aj_n}.
\]
Moreover,
\small
\begin{align*}
&\left(\rho_{\vec{c}}\left(\alpha_s-\frac{1}{p}\right)
+\delta-\alpha_sp\right)A \\
&\qquad \leq
\left(\rho_{\vec{c}}\left(\alpha_s-\frac{1}{p}\right)
+\delta-\alpha_sp\right)
\frac{h+2\varepsilon}{\alpha_s-\frac{1}{p}} \\
&\qquad \leq
\left(
\left(
\rho_{\vec{c}}^{(p),\ast}(h)-1
-\left(
4+
\frac{h\left(h+2\varepsilon+\frac{1}{p}\right)}
{\left(\alpha_0+\frac{1}{p}\right)(h+2\varepsilon)}
\right)\delta
-M\varepsilon p
\right)
\frac{\alpha_s-\frac{1}{p}}{h}
+1+\delta-\alpha_sp
\right)
\frac{h+2\varepsilon}{\alpha_s-\frac{1}{p}} \\
&\qquad \leq
\left(\rho_{\vec{c}}^{(p),\ast}(h)-1-4\delta-M\varepsilon p\right)
\left(1+\frac{2\varepsilon}{h}\right)
-(h+2\varepsilon)p \\
&\qquad =
\rho_{\vec{c}}^{(p),\ast}(h)-(hp+1)-4\delta-2\varepsilon p-M\varepsilon p
+\frac{2\varepsilon}{h}
\left(\rho_{\vec{c}}^{(p),\ast}(h)-1-4\delta-M\varepsilon p\right) \\
&\qquad \leq
\rho_{\vec{c}}^{(p),\ast}(h)-(hp+1)-4\delta-2\varepsilon p.
\end{align*}
\normalsize
Here, the first two estimates follow from Condition~\eqref{eq:delta} and
from the definition of $I$, which ensures that
\[
\frac{h+2\varepsilon}{\alpha_s-\frac{1}{p}}
\leq
\frac{h+2\varepsilon+\frac{1}{p}}
{\alpha_0+\frac{1}{p}}.
\]
The last inequality follows from Relation~\eqref{eq:epsilon_neg} together
with~\eqref{eq:delta}, since
\[
M\varepsilon p+\frac{Mhp}{2}+1
<0<
\rho_{\vec{c}}^{(p),\ast}(h)-4\delta.
\]
This yields the desired contradiction.
\end{proof}

\end{document}
